\documentclass[11pt]{article}
\usepackage[margin=1in]{geometry}
\usepackage{amsmath,amssymb}
\usepackage{booktabs}

\title{Shared4 Multiplex Objective}
\author{}
\date{}

\begin{document}
\maketitle

\section{Setup and Notation}

Let \(x_i\) denote prompt \(i\). For each prompt, the policy samples
\(M=4\) multiplex thinking traces
\[
  z_{i1}, z_{i2}, z_{i3}, z_{i4}.
\]
From each thinking trace \(z_{im}\), the policy samples \(K=4\) answer
continuations
\[
  y_{im1}, y_{im2}, y_{im3}, y_{im4}.
\]
Thus each prompt produces \(MK=16\) scored answers. Let
\[
  r_{imj} = R(x_i, z_{im}, y_{imj})
\]
be the scalar reward assigned to the \(j\)-th answer continuation from the
\(m\)-th thinking trace. Let \(T_{im}\) be the set of token positions belonging
to the multiplex thinking trace, and let \(S_{imj}\) be the set of token
positions belonging to the final answer continuation.

\section{Thinking-Trace Reward and Advantage}

The reward assigned to a thinking trace is the average reward of its four
answer continuations:
\[
  \bar r^T_{im}
  =
  \frac{1}{K}\sum_{j=1}^{K} r_{imj}.
\]
For each prompt, the thinking-trace baseline and scale are computed over the
four trace-level rewards:
\[
  \mu_i^T = \frac{1}{M}\sum_{m=1}^{M} \bar r^T_{im},
  \qquad
  (\sigma_i^T)^2
  =
  \frac{1}{M-1}\sum_{m=1}^{M}(\bar r^T_{im}-\mu_i^T)^2.
\]
The normalized thinking advantage is
\[
  A^T_{im}
  =
  \frac{\bar r^T_{im}-\mu_i^T}{\sigma_i^T+\epsilon}.
\]

\section{Answer Reward and Advantage}

The answer-level baseline and scale are computed across all sixteen answers for
the same prompt:
\[
  \mu_i^A = \frac{1}{MK}\sum_{m=1}^{M}\sum_{j=1}^{K} r_{imj},
  \qquad
  (\sigma_i^A)^2
  =
  \frac{1}{MK-1}
  \sum_{m=1}^{M}\sum_{j=1}^{K}(r_{imj}-\mu_i^A)^2.
\]
The normalized answer advantage is
\[
  A^A_{imj}
  =
  \frac{r_{imj}-\mu_i^A}{\sigma_i^A+\epsilon}.
\]

\section{PPO Token Loss}

For a sampled token or soft-thinking candidate \(a\) at position \(t\), define
the importance ratio
\[
  \rho_t(a)
  =
  \exp\left(
    \log \pi_\theta(a \mid h_t)
    -
    \log \pi_{\theta_{\mathrm{old}}}(a \mid h_t)
  \right),
\]
where \(h_t\) is the generation history at position \(t\). The clipped PPO loss
for advantage \(A\) is
\[
  \ell(\rho, A)
  =
  \max\left(
    -A\rho,\,
    -A\,\mathrm{clip}(\rho, 1-\eta_{\mathrm{low}}, 1+\eta_{\mathrm{high}})
  \right).
\]
The implementation also supports the existing dual-clip lower bound for
negative advantages, matching the local VERL PPO code path.

\section{Thinking Loss}

For multiplex thinking tokens, let \(\mathcal Q_{imt}\) be the top-\(q\)
candidate support used by the soft-thinking rollout at token position \(t\).
The thinking loss is
\[
  L_T(\theta)
  =
  \frac{1}{B}
  \sum_{i=1}^{B}
  \frac{1}{M}
  \sum_{m=1}^{M}
  \frac{1}{|T_{im}|}
  \sum_{t\in T_{im}}
  \sum_{a\in \mathcal Q_{imt}}
  \ell\left(\rho_t(a), A^T_{im}\right).
\]
In the code path for this experiment, the soft-thinking top-\(q\) candidate
slots are optimized for thinking tokens, while final answer tokens use the
single realized answer token.

\section{Answer Loss}

For final answer tokens, the loss uses the per-answer advantage:
\[
  L_A(\theta)
  =
  \frac{1}{B}
  \sum_{i=1}^{B}
  \frac{1}{MK}
  \sum_{m=1}^{M}
  \sum_{j=1}^{K}
  \frac{1}{|S_{imj}|}
  \sum_{t\in S_{imj}}
  \ell\left(\rho_t(y_{imj,t}), A^A_{imj}\right).
\]

\section{Objective Variants}

The joint shared4 objective optimizes both the shared thinking trace and the
final answers:
\[
  L_{\mathrm{joint}}(\theta)
  =
  \lambda_T L_T(\theta) + \lambda_A L_A(\theta).
\]
For the submitted run, \(\lambda_T=\lambda_A=1\).

The thinking-only variant is
\[
  L_{\mathrm{thinking}}(\theta) = L_T(\theta).
\]
The answer-only variant is
\[
  L_{\mathrm{answer}}(\theta) = L_A(\theta).
\]

\section{Code Mapping}

\begin{center}
\begin{tabular}{ll}
\toprule
Mathematical object & Implementation location \\
\midrule
\(A^T_{im}\), \(A^A_{imj}\) &
\texttt{verl/trainer/ppo/core\_algos.py::compute\_shared4\_outcome\_advantage} \\
Shared trace grouping metadata &
\texttt{verl/trainer/ppo/ray\_trainer.py::attach\_shared4\_rollout\_metadata} \\
\(L_{\mathrm{joint}}\) &
\texttt{shared4\_joint} policy loss \\
\(L_{\mathrm{thinking}}\) &
\texttt{shared4\_thinking\_only} policy loss \\
\(L_{\mathrm{answer}}\) &
\texttt{shared4\_answer\_only} policy loss \\
\bottomrule
\end{tabular}
\end{center}

\end{document}
